\chapter{Rao-Blackwellised SQMC for Football Matches}
\label{ch:model}
In this chapter, we develop a correlated Gaussian state-space model for football match outcomes. The observation likelihood depends on only the two playing teams, so the effective dimension is far smaller than the $2K$ -dimensional latent state. We exploit this sparsity through Rao-Blackwellisation, integrating out the non-playing teams analytically and running SQMC on the four playing-team coordinates, overcoming the high-dimensional limitations of SQMC.

% ADD IN CURSE OF DIMENSIONALITY FROM \citep{chopingerber2017sqmcintroduction}

\section{Background}
\label{sec:rb-sqmc-background}
Football is a high-dimensional problem with sparse pairwise observations. Let each team's latent state be a two-dimensional vector of attacking and defensive strengths, $x_t^k = (x_t^{k, \text{attack}}, x_t^{k, \text{defence}})$ for $k \in \{1, 2, \ldots K\}$. Then, the latent state of all $K$ teams is $x_t \in \mathbb{R}^{2K}$. However, each match involves only the latent states of the two playing teams at time $t$, $x_t^{\mathcal{O}_t} = (x_t^{(i_t)}, x_t^{(j_t)}) \in \mathbb{R}^4$, through the local observation model $G_t(y_t \mid x_t^{\mathcal{O}_t})$, which represents the outcome of a single football match. The general state-space model notation introduced in Section~\ref{sec:background} is inefficient for this problem: it represents the full joint distribution over all teams, even though each likelihood contribution depends on only four coordinates.

% The general state-space model notation and the filtering and smoothing objectives were introduced in Section~\ref{sec:background}; here we specialise the relevant structure to football. Let $x_t^k$ denote the latent state of team $k$ at matchtime $t$, and let $\mathcal{O}_t = \{i_t,j_t\}$ denote the two teams involved in match $t$. The observation likelihood therefore has the local form $G_t(y_t \mid x_t^{\mathcal{O}_t})$.

The factorial hidden Markov model (fHMM) was introduced by Ghahramani and Jordan \citep{ghahramanijordan1995factorialhmm,ghahramanijordan1997factorialhmm} to tackle problems with structures that can be factorised into individual components. A factorial state-space model (fSSM) is the corresponding state-space construction, extending from the general state-space-model defined in Equation~\eqref{eq:state-space-model}. Its state $x$ is naturally partitioned into a collection of components $k$, each of which evolves independently of one another. The observation likelihood $G_t$ may depend jointly on a single component or on a collection of components. Let $\mathcal{O}_t \subseteq \{1,\ldots,K\}$ denote the subset of components involved in the observation at time $t$, and write $x_t^{\mathcal{O}_t} = (x_t^k)_{k \in \mathcal{O}_t}$ for the corresponding collection of components.

\begin{equation}
p(x_{0:T}^{1:K}, y_{1:T}) = \left[ \prod_{k=1}^K m_0(x_0^k) \prod_{t=1}^T m_t(x_t^k \mid x_{t-1}^k) \right] \prod_{t=1}^T G_t(y_t \mid x_t^{\mathcal{O}_t}).
\end{equation}

The fSSM model makes the following assumptions: (1) the initial latent states are \text{a priori} independent (independent given any observations); (2) evolution $m_t$ of each component $k$ is independent; (3) given the latent states $x_t^{\mathcal{O}_t}$, the observations $y_t$ are conditionally independent of all other components $-\mathcal{O}_t$. This reduces the computational complexity of the joint distribution from $\mathcal{O}(N \cdot K)$ to $\mathcal{O}(N + K)$ terms, which is a significant improvement for large $K$.

\citet{duffieldpowerrimella2024factorialssm} apply the fSSM to online skill rating and sports. In this context, our components can either represent an individual or a team's skill level. Their inference maintains a factorial filtering distribution, which factorises as a product of per-component marginals,
\begin{equation}
\label{eq:factorial-filter}
p(x_t \mid y_{1:t}) = \prod_{k=1}^K p(x_t^k \mid y_{1:t}).
\end{equation}
Because the observation at time $t$ involves only the pair $\mathcal{O}_t = \{i_t, j_t\}$, the prediction step propagates each component independently,
\begin{equation}
p(x_t^k \mid y_{1:t-1}) = \int m_t(x_t^k \mid x_{t-1}^k) \, p(x_{t-1}^k \mid y_{1:t-1}) \, dx_{t-1}^k,
\end{equation}
and the update step assimilates the match result in the joint distribution of the playing pair,
\begin{equation}
p(x_t^{\mathcal{O}_t} \mid y_{1:t}) \propto G_t(y_t \mid x_t^{\mathcal{O}_t}) \prod_{k \in \mathcal{O}_t} p(x_t^k \mid y_{1:t-1}).
\end{equation}
For each component $k \in \mathcal{O}_t$, the marginal filtering distribution is then recovered by integrating the joint update over the remaining components of the playing pair,

\begin{equation}
p(x_t^k \mid y_{1:t}) = \int p(x_t^{\mathcal{O}_t} \mid y_{1:t}) \, dx_t^{\mathcal{O}_t \setminus \{k\}}, \qquad k \in \mathcal{O}_t,
\end{equation}

and we obtain the filtering distribution in the factorial form of Equation~\eqref{eq:factorial-filter}. Figure~\ref{fig:fssm-independence} contrasts the conditional independence structure of a general state-space model with that of a fSSM: in the latter, each component evolves independently and the observation at time $t$ couples only the components in $\mathcal{O}_t$. The main approximation is the repeated projection of the coupled posterior back to a product of component-level marginals, which discards the cross-component covariance after every match. This approximation is computationally attractive because it avoids representing the full joint distribution over all components, which contains $\mathcal{O}(N \cdot K)$ terms, while this new representation only contains $\mathcal{O}(N + K)$ terms.

\begin{figure}[htbp]
\centering
\includegraphics[width=0.9\textwidth]{m_qlae035f1.jpeg}
\caption{Conditional independence structure of a general state-space model (left) and of a factorial state-space model with a pairwise observation model (right), reproduced from \citet{duffieldpowerrimella2024factorialssm}. In the fSSM, each component evolves independently and the observation at time $t$ couples only the components in $\mathcal{O}_t$.}
\label{fig:fssm-independence}
\end{figure}

Although efficient, fSSM is a restrictive model for football. Shared players, clubs, leagues and tournaments are potential sources of dependence between national teams that are not represented by a model whose team processes are \textit{a priori} independent (i.e. assumed to be independent without any observations). Furthermore, match outcomes create dependencies that inform us on the strengths and uncertainty of other teams (e.g. if team $A$ beats team $B$, and team $B$ beats team $C$, we can infer that team $A$ is likely stronger than team $C$). This information is discarded in the fSSM by approximating the filtering distribution as a product of independent team marginals as defined in Equation~\eqref{eq:factorial-filter}. 

% The resulting marginals $p(x_t^k \mid y_{1:t})$ are still useful for prediction, but they do not retain the joint uncertainty induced by the observed competition history.

The model introduced in this chapter uses the original state-space model defined in Equation~\eqref{eq:state-space-model}. The full latent state is high-dimensional, so a direct particle filter costs $\mathcal{O}(T \cdot N \cdot K)$ per run, where $N$ is the number of particles and $T$ the number of observations. To avoid this, we restrict the transition model to linear-Gaussian and apply Rao-Blackwellisation: the non-playing teams are integrated out analytically, so that only the four coordinates of the playing pair are sampled. We then run SMC or SQMC on this $(2 \times 2)$-dimensional block and recover the posterior via the Kalman filter, reducing the cost to $\mathcal{O}(N \cdot (2 \times 2))$, independent of $K$.

\section{Model}
\label{sec:model}
We represent the latent states of each team as a 2-dimensional vector $x = (x^{\text{attack}}, x^{\text{defence}})$ and the full latent state of $K$ teams is $x \in \mathbb{R}^{2K}$ as mentioned in the previous Section~\ref{sec:rb-sqmc-background}. Our initial distribution of latent states is given by 

\begin{equation}
    x_0 \in \mathbb{R}^{2K}, \qquad x_0 \sim \mathcal{N}_{2K}(\mu_0, \Sigma_0)
\end{equation}

We make the further assumption that the relationship between attack and defence of all teams is the same, represented by the matrix $B \in \mathbb{R}^{2 \times 2}$. This assumption allows us to simplify the covariance matrix into a Kronecker product of the form $\Sigma_0 = \Gamma_0 \otimes B$, where $\Gamma_0 \in \mathbb{R}^{K \times K}$ is a covariance matrix that represents the relationship between teams. We require the conditions $\Gamma_0 \succ 0$, $B \succ 0$ and $\det(B) = 1$. The condition $\det(B) = 1$ ensures that the scale of the Kronecker product is identifiable since $(c\Gamma_0) \otimes (B/c) = \Gamma_0 \otimes B$ (i.e. we fix the determinant of $B$ so that $\Gamma_0$ is the free parameter in the Kronecker product). This paramterization is a natural choice for football as we can expect that attack and defence strengths of teams across all teams should share a similar relationship. It also reduces the number of free parameters in the initial covariance matrix from $K(2K+1)$ to $K(K+1)/2 + 2$ as seen in Proposition~\ref{prop:kronecker-dim}, allowing us to efficiently compute the models in the filtering, smoothing and parameter estimation steps. 

% Proposition~\ref{prop:kronecker-dim} in Appendix~\ref{app:kronecker-proof} shows that this parameterisation reduces the number of free parameters of the initial covariance from $K(2K+1)$ to $K(K+1)/2 + 2$; the reduction is a factor of approximately four, but the parameter count remains of order $\mathcal{O}(K^2)$. This also allows us to efficiently compute the models in the following Section~\ref{sec:rb-smc} and \ref{sec:rb-sqmc}. 

We model the evolution of latent states as an Ornstein-Uhlenbeck process, a mean-reverting stochastic process, given by

\begin{equation}
    x_t \mid x_{t-1} \sim \mathcal{N}(\mu_0 + \phi_t (x_{t-1} - \mu_0), Q_t)
\end{equation}

where we define $\phi_t = \exp(-\kappa \Delta t_t)$, $\Delta t_t = t_t - t_{t-1}$ and $Q_t = (1 - \phi_t^2)\Sigma_0$ is the covariance matrix of the process noise. $\kappa$ is the mean-reversion rate and $\Delta t_t$ is the time difference between $t$ and $t-1$, measured in days. This behaviour is desirable in football as we expect average skill levels to revert to a long-run marginal distribution of $\mathcal{N}(\mu_0, \Sigma_0)$ over time without any new observations.

As mentioned before, each match only involves two teams, so we represent the latent states of the two playing teams as $x_t^{\mathcal{O}_t} = (x_t^{(i_t)}, x_t^{(j_t)})$. Our observation is a bivariate vector of scores for team $i_t$ and team $j_t$, represented as $y_t = (y_t^{(i_t)}, y_t^{(j_t)})$. The likelihood function follows the bivariate Poisson distribution that \citet{karlisntzoufras2003bivariatepoisson} proposed for sports data.

\begin{equation}
\label{eq:likelihood}
    G_t(y_t \mid x_t^{\mathcal{O}_t}) = e^{-(\lambda_1 + \lambda_2 + \lambda_3)} \frac{\lambda_1^{y_t^{(i_t)}}}{y_t^{(i_t)}!} \frac{\lambda_2^{y_t^{(j_t)}}}{y_t^{(j_t)}!} \sum_{k=0}^{\min(y_t^{(i_t)}, y_t^{(j_t)})} \binom{y_t^{(i_t)}}{k} \binom{y_t^{(j_t)}}{k} k! \left( \frac{\lambda_{3}}{\lambda_1 \lambda_2} \right)^k
\end{equation}

where $\lambda_1 = \exp(\alpha + x_t^{\text{att}, i_t} - x_t^{\text{def}, j_t})$, $\lambda_2 = \exp(\alpha + x_t^{\text{att}, j_t} - x_t^{\text{def}, i_t})$, $\lambda_3 = \exp(\beta)$. The parameter $\alpha$ and $\beta$ are static parameters of the model. $\alpha$ can be used to model home-team advantage by estimating $\alpha^h$ in $\lambda_1$ and $\alpha^a$ in $\lambda_2$. One could also investigate the effect of the home advantage by testing the hypothesis $\alpha^h > \alpha^a$.

\section{Rao-Blackwellised SMC (RB-SMC)}
\label{sec:rb-smc}
In our previous Section~\ref{sec:model}, we state that only the latent states of the two playing teams $x_t^{\mathcal{O}_t}$ are relevant to the likelihood function $G_t$. Defining the time-varying non-playing latent states as $x_t^{\mathcal{R}_t} = x_t^{-\mathcal{O}_t}$, and interpreting $x_{0:T}^{\mathcal{O}}$, the Rao-Blackwellised SMC (RB-SMC) algorithm therefore targets the marginal posterior of the coordinates of the playing teams $p(x_{0:T}^{\mathcal{O}} \mid y_{1:T})$ and recovers the non-playing teams $p(x_{0:T}^{\mathcal{R}} \mid x_{0:T}^{\mathcal{O}}, y_{1:T})$ in closed form. This is similar to the Rao-Blackwellised particle filter construction of \citet{doucetdefreitasmurphyrussell2000rbpf}, applied to the sparse football likelihood. We can then factorize the full joint posterior as follows:

% In Section~\ref{sec:smc}, we described the SMC algorithm for a general state-space model. 

% The RB-SMC targets the marginal posterior of the sequence of sampled coordinates and recovers the non-playing teams in closed form instead. 

% In Section~\ref{sec:smc}, we described the SMC algorithm for a general state-space model. 

\begin{equation}
    p(x_{0:T} \mid y_{1:T}) = p(x_{0:T}^{\mathcal{O}} \mid y_{1:T}) \cdot p(x_{0:T}^{\mathcal{R}} \mid x_{0:T}^{\mathcal{O}}, y_{1:T})
\end{equation}

Since the likelihood depends only on $x_{0:T}^{\mathcal{O}}$, the second factor simplifies to
\begin{equation}
    p(x_{0:T}^{\mathcal{R}} \mid x_{0:T}^{\mathcal{O}}, y_{1:T}) = p(x_{0:T}^{\mathcal{R}} \mid x_{0:T}^{\mathcal{O}}).
\end{equation}

As stated before, our transition model is linear-Gaussian, so the covariance recursions are deterministic and do not depend on the sampled coordinates or the particle paths; hence $\Sigma_t$ is shared across all particles while the means $\mu_t^{(i)}$ carry the particle-specific information. The RB-SMC algorithm has the same steps as the SMC algorithm as defined in Algorithm~\ref{alg:pf}, the only difference is that we only sample the latent state of playing teams $x_t^{\mathcal{O}_t}$ at each time step and compute the marginal posterior for the non-playing teams in closed form. The RB-SMC algorithm is as follows:

\textbf{1 Prediction}: Compute the predictive mean and covariance matrix for the full state:

\begin{equation}
    \mu_{t|t-1}^{(i)} = \mu_0 + \phi_t (\mu_{t-1}^{(i)} - \mu_0) \nonumber
\end{equation}

\begin{equation}
    \Sigma_{t|t-1} = \phi_t^2 \Sigma_{t-1} + (1 - \phi_t^2) \Sigma_0 \nonumber
\end{equation}

\textbf{2 Bootstrap Proposal Sampling}: For particle $i$, the full predictive distribution is
\begin{equation}
x_t \mid I_{t-1}=i,y_{1:t-1} \sim \mathcal{N}\left(\mu_{t\mid t-1}^{(i)}, \Sigma_{t\mid t-1} \right). \end{equation}

At each observation, only the two playing teams are sampled. We extract the predictive mean and covariance for the playing teams' coordinates $x_t^{\mathcal{O}_t}$, and the following predictive marginal under the bootstrap proposal is defined as,

\begin{equation}
    q_t^{(i)}(x_t^{\mathcal{O}_t}) = p(x_t^{\mathcal{O}_t} \mid I_{t-1} = i, y_{1:t-1}) = \mathcal{N}(\mu_{t|t-1}^{\mathcal{O}_t,(i)}, \Sigma_{t|t-1}^{\mathcal{O}_t\mathcal{O}_t}) \nonumber
\end{equation}

\textbf{3 Computing log-weights}: For the propagate-weight-resample ordering, compute the unnormalised log-weights for each particle $(i)$ using the previous weight and the likelihood function $G_t$,

\begin{equation}
    \log \tilde{w}_t^{(i)} = \log w_{t-1}^{(i)} +  \log G_t(y_t \mid x_t^{\mathcal{O}_t, (i)}) \nonumber
\end{equation}

and normalise,
\begin{equation}
    w_t^{(i)} = \frac{\tilde{w}_t^{(i)}}{\sum_{n=1}^N \tilde{w}_t^{(n)}}. \nonumber
\end{equation}

\textbf{4 Exact Marginalisation of Non-Playing Teams}: Given the predictive joint distribution and conditioning on the playing teams' coordinates $x_t^{\mathcal{O}_t}$, the non-playing coordinates $x_t^{\mathcal{R}_t}$ are Gaussian with a closed-form mean and covariance matrix defined by the Kalman filter equations,

\begin{equation}
    x^{\mathcal{R}_t,(i)}_{t \mid t-1} \mid x^{\mathcal{O}_t,(i)}_{t} \sim \mathcal{N}\left( \mu^{\mathcal{R}|\mathcal{O},(i)}_{t \mid t-1}, \Sigma^{\mathcal{R}|\mathcal{O}}_{t \mid t-1} \right) \nonumber
\end{equation}

\begin{align}
    \mu^{\mathcal{R}_t|\mathcal{O}_t,(i)}_{t \mid t-1} &= \mu_{t|t-1}^{\mathcal{R}_t} + K_t \left( x_t^{\mathcal{O}_t,(i)} - \mu_{t|t-1}^{\mathcal{O}_t} \right), \nonumber \\
    \Sigma^{\mathcal{R}_t|\mathcal{O}_t}_{t \mid t-1} &= \Sigma_{t|t-1}^{\mathcal{R}_t\mathcal{R}_t} - K_t \Sigma_{t|t-1}^{\mathcal{O}_t\mathcal{R}_t}, \nonumber \\
    K_t &= \Sigma_{t|t-1}^{\mathcal{R}_t\mathcal{O}_t} \left( \Sigma_{t|t-1}^{\mathcal{O}_t\mathcal{O}_t} \right)^{-1}. \nonumber
\end{align}

In Section~\ref{sec:smc}, we made the assumption that our initial covariance is a Kronecker product of the form $\Sigma_0 = \Gamma_0 \otimes B$. Since our transition model is linear-Gaussian, the Kronecker structure is preserved recursively by induction, 

\begin{align}
    \Sigma_{t \mid t-1} &= \phi_t^2 \Sigma_{t-1} + (1 - \phi_t^2) \Sigma_0 \nonumber \\
    &= \phi_t^2 (\Gamma_{t-1} \otimes B) + (1 - \phi_t^2) (\Gamma_0 \otimes B) \nonumber \\
    &= \left( \phi_t^2 \Gamma_{t-1} + (1 - \phi_t^2) \Gamma_0 \right) \otimes B \nonumber \\
    &= \Gamma_{t|t-1} \otimes B. \nonumber
\end{align}

so that the covariance blocks for the playing and non-playing teams factor as:

\begin{equation}
    \Sigma_{t|t-1}^{\mathcal{O}_t\mathcal{O}_t} = \Gamma^{\mathcal{O}_t\mathcal{O}_t}_{t|t-1} \otimes B, \qquad
    \Sigma_{t|t-1}^{\mathcal{R}_t\mathcal{R}_t} = \Gamma^{\mathcal{R}_t\mathcal{R}_t}_{t|t-1} \otimes B, \qquad
    \Sigma_{t|t-1}^{\mathcal{O}_t\mathcal{R}_t} = \Gamma^{\mathcal{O}_t\mathcal{R}_t}_{t|t-1} \otimes B \nonumber
\end{equation}

where $\Gamma^{\mathcal{O}_t\mathcal{O}_t}_{t|t-1} \in \mathbb{R}^{2 \times 2}$ is the submatrix of $\Gamma_{t|t-1}$ for the two playing teams and $\Gamma^{\mathcal{R}_t\mathcal{R}_t}_{t|t-1} \in \mathbb{R}^{(K-2) \times (K-2)}$ is the submatrix for the non-playing teams. Using the mixed-product property $(A \otimes B)(C \otimes D) = (AC) \otimes (BD)$ and the inverse property $(A \otimes B)^{-1} = A^{-1} \otimes B^{-1}$, the Kalman gain factorises as:

\begin{align}
    K_t &= \Sigma_{t|t-1}^{\mathcal{R}_t\mathcal{O}_t} \left( \Sigma_{t|t-1}^{\mathcal{O}_t\mathcal{O}_t} \right)^{-1} \nonumber \\
    &= \left( \Gamma^{\mathcal{R}_t\mathcal{O}_t}_{t|t-1} \otimes B \right) \left( \Gamma^{\mathcal{O}_t\mathcal{O}_t}_{t|t-1} \otimes B \right)^{-1} \nonumber \\
    &= \left( \Gamma^{\mathcal{R}_t\mathcal{O}_t}_{t|t-1} \otimes B \right) \left( (\Gamma^{\mathcal{O}_t\mathcal{O}_t}_{t|t-1})^{-1} \otimes B^{-1} \right) \nonumber \\
    &= \left( \Gamma^{\mathcal{R}_t\mathcal{O}_t}_{t|t-1} (\Gamma^{\mathcal{O}_t\mathcal{O}_t}_{t|t-1})^{-1} \right) \otimes \left( B B^{-1} \right) \nonumber \\
    &= \tilde{K}_t \otimes I_2 \nonumber
\end{align}

\begin{equation}
    \tilde{K}_t = \Gamma^{\mathcal{R}_t\mathcal{O}_t}_{t|t-1} (\Gamma^{\mathcal{O}_t\mathcal{O}_t}_{t|t-1})^{-1} \nonumber
\end{equation}

and the conditional covariance of the non-playing teams becomes

\begin{equation}
    \Sigma_{t|t-1}^{\mathcal{R}_t|\mathcal{O}_t} = \left( \Gamma^{\mathcal{R}_t\mathcal{R}_t}_{t|t-1} - \tilde{K}_t \Gamma^{\mathcal{O}_t\mathcal{R}_t}_{t|t-1} \right) \otimes B \nonumber
\end{equation}

To preserve the Kronecker structure after conditioning, we update the full component mean and the full covariance. Define the team-space Kalman gain

\begin{equation}
    \overline{K}_t = \Gamma_{t|t-1}^{:\mathcal{O}_t} \left( \Gamma_{t|t-1}^{\mathcal{O}_t\mathcal{O}_t} \right)^{-1}, \nonumber
\end{equation}

where $\Gamma_{t|t-1}^{:\mathcal{O}_t}$ contains all rows and the two playing-team columns of $\Gamma_{t|t-1}$. The complete component mean is updated as

\begin{equation}
    \mu_t^{(i)} = \mu_{t|t-1}^{(i)} + (\overline{K}_t \otimes I_2) \left( x_t^{\mathcal{O}_t,(i)} - \mu_{t|t-1}^{\mathcal{O}_t,(i)} \right). \nonumber
\end{equation}

The rows of $\overline{K}_t$ belonging to $\mathcal{O}_t$ form the identity, so the playing-team entries of $\mu_t^{(i)}$ equal the sampled entries $x_t^{\mathcal{O}_t,(i)}$, while the $\mathcal{R}_t$ rows recover the conditional mean of the non-playing teams in closed form. The team-space filtered covariance $\Gamma_t$ is defined by the Schur complement

\begin{equation}
    \Gamma_t = \Gamma_{t|t-1} - \overline{K}_t \Gamma_{t|t-1}^{\mathcal{O}_t} = \Gamma_{t|t-1} - \Gamma_{t|t-1}^{:\mathcal{O}_t} \left( \Gamma_{t|t-1}^{\mathcal{O}_t\mathcal{O}_t} \right)^{-1} \Gamma_{t|t-1}^{\mathcal{O}_t} \nonumber
\end{equation}

and hence to recover the full filtered covariance $\Sigma_t$, one can simply compute the Kronecker product of $\Gamma_t$ with $B$. Because the rows and columns of $\Gamma_t$ corresponding to the playing teams $\mathcal{O}_t$ are zero, the full filtered covariance is sparse.

\begin{equation}
    \Sigma_t = \Gamma_t \otimes B = 
    \begin{pmatrix}
        \mathbf{0}_{4 \times 4} & \mathbf{0}_{4 \times 2(K-2)} \\
        \mathbf{0}_{2(K-2) \times 4} & \Gamma_t^{\mathcal{R}_t\mathcal{R}_t} \otimes B
    \end{pmatrix},
    \nonumber
\end{equation}

where $\mathbf{0}_{4 \times 4}$ is the zero block for the four coordinates of the two playing teams, whose latent states are now known (sampled) values with no remaining uncertainty, and $\Gamma_t^{\mathcal{R}_t\mathcal{R}_t} \otimes B$ is the non-zero conditional covariance of the non-playing teams. Intuitively, once a team has played a match, its skill level is exactly known at that point in time. Additionally, as mentioned in Section~\ref{sec:model}, the Kronecker structure reduces the number of free parameters in the covariance matrix from $K(2K+1)$ to $K(K+1)/2 + 2$. So then computing $\Gamma_t$ at each time step costs $\mathcal{O}(K^3)$ instead of $\mathcal{O}((2K)^3)$, reducing the constant factor by 8.  

% This full update encodes both effects of conditioning in a single matrix: the rows and columns of $\Gamma_t$ corresponding to the playing teams $\mathcal{O}_t$ are zero, because their coordinates are now known (sampled) values with no remaining uncertainty, while the $\mathcal{R}_t\mathcal{R}_t$ block of $\Gamma_t$ is the non-zero conditional covariance of the non-playing teams, which is precisely the Schur complement $\Sigma_{t|t-1}^{\mathcal{R}_t|\mathcal{O}_t}$ derived above.

% The Kronecker structure provides computational benefits for the algorithm. Instead of computing Kalman recursions on the full $2K \times 2K$ covariance matrix $\Sigma$, our Kalman recursions are computed on the $K \times K$ between-team covariance matrix $\Gamma$ and the $2 \times 2$ matrix $B$. The covariance recursion is shared across all particles, so we only need to compute it for $T$ time steps, rather than $N \times T$ times for all particles. With dense storage, updating the covariance costs $\mathcal{O}(K^2)$ per observation after a $2 \times 2$ solve, while updating all particle means costs $\mathcal{O}(NK)$. Thus the method reduces constants and avoids particle-specific covariance recursions.

\textbf{5 Resampling}: Resample the particles according to the normalised weights $\{w_t^{(i)}\}_{i=1}^N$ and set the weights to $w_t^{(i)} = 1/N$ for all particles.

\section{RB-SMC Smoothing}
\label{sec:rb-smc-smoothing}

In standard particle-based Forward Filtering Backward Simulation (FFBSi), the backward state at time $t$ is sampled from the discrete set of filtering particles, with probabilities determined by their filtering weights and compatibility with the sampled state at time $t+1$ \citep{godsilldoucetwest2004montecarlosmoothing}. 

% In our RB-SMC formulation we cannot perform backward simulation by drawing samples one step at a time from the filter results (as in the standard FFBSi algorithm), because the complementary coordinates $x_t^{\mathcal{R}}$ are marginalized analytically and their conditional distribution depends on the entire observed-coordinate trajectory $x_{1:t}^{\mathcal{O}}$, not just $x_t^{\mathcal{O}}$. Thus we need to sample the whole trajectories from the sample histories produced by the RB-SMC algorithm, similar to the case in \citep{sarkkabunchgodsill2012rbpsmoother}.

In the RB-SMC formulation, discrete point-particle FFBSi cannot be applied directly, because each filtering particle represents a Gaussian component rather than a complete state. Using the fact that our transition is linear and Gaussian, our terminal posterior filtering distribution can be approximated as a weighted Gaussian mixture with a shared covariance matrix $\Sigma_t$ across all particles.

\begin{equation}
    p(x_{T} \mid y_{1:T}) \approx \sum_{i=1}^N w_{T}^{(i)} \mathcal{N}(\mu_T^{(i)}, \Sigma_T)
\end{equation}

This means that a one-step backward kernel is available in closed form to obtain a smoothed trajectory from terminal time $T$ to start time $0$: first select a Gaussian component using its backward weight, and then draw the full state from that component's conditional Gaussian distribution. 

\subsection{Deriving the one-step backward kernel}

The one-step backward kernel satisfies
\begin{equation}
p(x_t \mid x_{t+1}, y_{1:t}) \propto p(x_t \mid y_{1:t})\, p(x_{t+1} \mid x_t).
\end{equation}
To derive this kernel, apply Bayes' rule to the filtering mixture at time $t$,
\begin{equation}
p(x_t \mid y_{1:t}) \approx \sum_{i=1}^N w_t^{(i)} \mathcal{N}(x_t \mid \mu_t^{(i)}, \Sigma_t),
\end{equation}
and the linear-Gaussian transition $p(x_{t+1} \mid x_t) = \mathcal{N}(x_{t+1} \mid a_{t+1}(x_t), Q_{t+1})$, where $a_{t+1}(x_t) = \mu_0 + \phi_{t+1}(x_t - \mu_0)$ and $Q_{t+1} = (1 - \phi_{t+1}^2)\Sigma_0$. Substituting the mixture and the transition,
\begin{equation}
p(x_t \mid x_{t+1}, y_{1:t}) \propto \sum_{i=1}^N w_t^{(i)} \mathcal{N}(x_t \mid \mu_t^{(i)}, \Sigma_t)\, \mathcal{N}(x_{t+1} \mid a_{t+1}^{(i)}, R_{t+1}),
\end{equation}
where $a_{t+1}^{(i)} = \mu_0 + \phi_{t+1}(\mu_t^{(i)} - \mu_0)$ is the predictive mean of component $i$ and $R_{t+1} = \phi_{t+1}^2 \Sigma_t + (1 - \phi_{t+1}^2)\Sigma_0$ is the corresponding predictive covariance. Because the second Gaussian factor does not depend on $x_t$, it acts as a component weight, so the backward kernel is itself a Gaussian mixture,
\begin{equation}
p(x_t \mid x_{t+1}, y_{1:t}) \approx \sum_{i=1}^N w_{t \mid t+1}^{(i)} \mathcal{N}(x_t \mid \mu_{t \mid t+1}^{(i)}, \Sigma_{t \mid t+1}),
\end{equation}
with backward weights
\begin{equation}
w_{t \mid t+1}^{(i)} \propto w_t^{(i)} \mathcal{N}(x_{t+1} \mid a_{t+1}^{(i)}, R_{t+1}),
\end{equation}
and, by the Rauch-Tung-Striebel (RTS) equations, the corrected mean and covariance of each component
\begin{equation}
\mu_{t \mid t+1}^{(i)} = \mu_t^{(i)} + J_t \left( x_{t+1} - a_{t+1}^{(i)} \right), \qquad \Sigma_{t \mid t+1} = \Sigma_t - J_t R_{t+1} J_t^T, \qquad J_t = \Sigma_t \phi_{t+1} R_{t+1}^{-1}.
\end{equation}
Sampling from this mixture therefore proceeds in two steps: draw the component index $I_t \sim \operatorname{Categorical}(w_{t \mid t+1}^{1:N})$, then draw the full state $x_t \sim \mathcal{N}(\mu_{t \mid t+1}^{(I_t)}, \Sigma_{t \mid t+1})$. Repeating this procedure from $t = T-1$ down to $t = 0$ produces a single complete smoothed trajectory $\{X_0^*, \dots, X_T^*\}$.

\subsection{The RB-SMC smoothing algorithm}

The RB-SMC smoothing algorithm is as follows:

\textbf{1 Initialisation} ($t = T$): Sample the terminal particle index $I_T$ using the normalised weights $w_T^{(i)}$, and draw the full terminal state from the corresponding Gaussian mixture,
\begin{equation}
    I_T \sim \operatorname{Categorical}(w_T^{1:N}), \qquad X_T^* \sim \mathcal{N}(\mu_T^{(I_T)}, \Sigma_T). \nonumber
\end{equation}

\textbf{2 Backward recursion} (for $t = T-1, T-2, \dots, 0$): Given the already-sampled state $X_{t+1}^*$, the backward kernel derived in Section~\ref{sec:rb-smc-smoothing} is a Gaussian mixture, and the recursion proceeds in three sub-steps.

\textbf{2a Compute the backward weights}: For every particle $i$, propagate the filtering mean through the OU dynamics and form the predictive covariance,
\begin{equation}
    a_{t+1}^{(i)} = \mu_0 + \phi_{t+1}(\mu_t^{(i)} - \mu_0), \qquad R_{t+1} = \phi_{t+1}^2 \Sigma_t + (1 - \phi_{t+1}^2) \Sigma_0, \nonumber
\end{equation}
then weight each component by its filtering weight and the predictive density of the conditioning state $X_{t+1}^*$,
\begin{equation}
    w_{t \mid t+1}^{(i)} \propto w_t^{(i)} \mathcal{N}\left( X_{t+1}^* \mid a_{t+1}^{(i)}, R_{t+1} \right), \nonumber
\end{equation}
normalised so that $\sum_{i=1}^N w_{t \mid t+1}^{(i)} = 1$.

\textbf{2b Select and correct the backward component}: Sample the particle index from the backward weights,
\begin{equation}
    I_t \sim \operatorname{Categorical}(w_{t \mid t+1}^{1:N}), \nonumber
\end{equation}
and condition the chosen component on $X_{t+1}^*$ through the RTS equations with the gain
\begin{equation}
    J_t = \Sigma_t \phi_{t+1} R_{t+1}^{-1}, \nonumber
\end{equation}
\begin{align}
    \mu_{t \mid t+1}^{(I_t)} &= \mu_t^{(I_t)} + J_t \left( X_{t+1}^* - a_{t+1}^{(I_t)} \right), \nonumber \\
    \Sigma_{t \mid t+1} &= \Sigma_t - J_t R_{t+1} J_t^T. \nonumber
\end{align}

\textbf{2c Draw the smoothed state}: Sample the full state from the corrected component,
\begin{equation}
    X_t^* \sim \mathcal{N}\left( \mu_{t \mid t+1}^{(I_t)}, \Sigma_{t \mid t+1} \right). \nonumber
\end{equation}

The recursion terminates at $t = 0$, yielding the smoothed trajectory $\{X_0^*, \dots, X_T^*\}$. Repeating this procedure $M$ times produces $M$ independent smoothed trajectories $\{X_{0:T}^{*(m)}\}_{m=1}^M$. The resulting algorithm is similar to the pure particle smoother of \citet{godsilldoucetwest2004montecarlosmoothing}, but the backward weights combine the forward filtering weights with the predictive density, while the RTS backward estimates determine the Gaussian mixture component from which the full state is drawn \citep{sarkkabunchgodsill2012rbpsmoother}. This backward-simulation construction for a Rao-Blackwellised particle filter on a conditionally linear-Gaussian model follows \citet{lindstenschonbunchgodsill2012rbpsmoother}, and the use of backward simulation in particle methods more generally is discussed by \citet{lindstenschon2011backwardsimulation}. This process is extremely parallelisable due to the independence of each smoothed trajectory.

\section{Rao-Blackwellised SQMC (RB-SQMC)}
\label{sec:rb-sqmc}

Since the likelihood function depends only on the latent states of the two playing teams, our 'effective dimension' is $d = 2 \times 2 = 4$ instead of the full state dimension of $2K$. This motivates the use of SQMC methods, which are known to perform better than SMC methods in low-dimensional settings as seen in Section~\ref{sec:sqmc-properties}. 

% We now extend the RB-SMC algorithm to use low-discrepancy sampling for the four-dimensional playing-team coordinates since QMC methods are shown to have better properties than standard SMC methods in Section~\ref{sec:sqmc-properties}.

\citep{chopingerber2017sqmcintroduction} show that SQMC methods can be applied to high-dimensional state-space models if the transition and likelihood functions only depend on a subset of the state coordinates. In our case, we can Hilbert-sort the $2K$-dimensional particles using the $4$-dimensional predicted means of the two playing teams, $\Lambda_t(\mu_{t \mid t-1}^{\mathcal{O}_t, (i)})$, with the assumption that the $4$-dimensional predicted means capture the particle diversity in the coordinates that determine the current likelihood. The likelihood is local because $G_t(x_{t-1}, x_t) = G_t(x_t^{\mathcal{O}_t})$, so the potential function depends only on the propagated coordinates of the two playing teams. These two characteristics means that our conditions are satisfied for preserving the low-discrepancy properties of the QMC point set, as shown in \citet{chopingerber2017sqmcintroduction}. 

Similarly in Section~\ref{sec:rb-smc}, the RB-SQMC is the same as SQMC algorithm defined in Algorithm~\ref{alg:sqmc}, with the only changes being made in the indexing of the teams being played $x_t^{\mathcal{O}}$. Our algorithm, the RB-SQMC, follows these steps:

% Unlike the SQMC algorithm described in Section~\ref{sec:sqmc},

% In our model, the likelihood at each observation only depends on the attack and defence strengths of the two playing teams. Rao-Blackwellisation analytically marginalises the particles of the non-playing teams, so only four coordinates are sampled.



% We Hilbert sort the full $2K$-dimensional particles using the four-dimensional predicted means of the two playing teams, $\Lambda_t(\mu_{t \mid t-1}^{\mathcal{O}_t, (i)})$. Once an ancestor is selected, its complete components are retained and propagated via the OU process. This is analogous to step 1 in Section~\ref{sec:rb-smc}. The likelihood is local because $G_t(x_{t-1}, x_t) = G_t(x_t^{\mathcal{O}_t})$, so the potential function depends only on the propagated coordinates of the two playing teams, matching the locality assumption in Section~\ref{sec:model}; this is the setting in which Chopin and Gerber (2017) show that sorting on a projected coordinate can preserve the QMC structure. 

% The four-dimensional sorting key $\Lambda(\mu_{t-1}^{(i)})$ captures the coordinates on which the current likelihood depends. Because the full Rao-Blackwellised component also carries non-playing means that influence future transitions, we retain the complete component after selection, so the projected ordering is a dimension-reduction heuristic rather than a full-state ordering, and we do not claim a formal convergence guarantee for it beyond the locality setting above.

% In our model, the likelihood at each observation depends only on the attacking and defensive strengths of the two teams involved in the match. Rao-Blackwellisation allows the remaining coordinates to be marginalized analytically, so only four latent coordinates are required to be propagated via sampling. This enable us to leverage the benefits of SQMC in the SMC operation.

% We Hilbert-sort the Rao-Blackwellised components using the four-dimensional state of the teams involved in the current match. The complete ancestor component, including the non-playing team means, is nevertheless retained and propagated. This projected ordering does not satisfy the exact transition-sufficiency condition of Chopin and Gerber (2017), so we treat it as a dimension-reduction heuristic.

% Our likelihood function satisfies the condition where $G_t$ depends only on $X_t$ and $\Lambda(x_{t - 1})$, where $\Lambda(x_{t - 1}) = x^{\mathcal{O},(i)}_{t - 1}$.

% We approximate the conditional distribution of non-playing teams as depending on the current playing teams state. This does not satisfy the exact transition-sufficiency condition used by Chopin and Gerber (2017), because non-playing component means remain relevant to later transitions. It is therefore treated as a dimension-reduction heuristic.

% The effective propagation dimension is therefore $d= 2 \times 2 = 4$, rather than the full state dimension $2K$; the full covariance matrix has dimension $2K \times 2K$. Including the additional coordinate used for ancestor selection, each SQMC point has dimension $d+1=5$. Our algorithm, the RB-SQMC, follows these steps,

% This dimension-reduction strategy follows the application described in Section~3.1 of \citet{chopingerber2017sqmcintroduction} and helps preserve the low-discrepancy properties of the QMC point set. 

\textbf{1 Generate N QMC points}: Generate $N$ QMC or RQMC points $(u_t^{1:N}, v_t^{1:N})$ of dimension $1 + (2 \times 2) = 5$, where $u_t$ is the first component, used for ancestor selection, and $v_t$ is the vector of $d = 4$ components used to propagate the latent states of the two playing teams, $x_t^i$ and $x_t^j$.

\textbf{2 Hilbert Sort}: Apply the deterministic OU prediction to each Rao-Blackwellised component. The predictive mean and covariance of the full state are
\begin{equation}
    \mu_{t\mid t-1}^{(i)} = \mu_0 + \phi_t \left( \mu_{t-1}^{(i)} - \mu_0 \right), \qquad
    \Sigma_{t\mid t-1} = \phi_t^2 \Sigma_{t-1} + \left(1 - \phi_t^2\right) \Sigma_0, \nonumber
\end{equation}
so that the playing-team block $\mu_{t\mid t-1}^{\mathcal{O}_t,(i)}$ and $\Sigma_{t\mid t-1}^{\mathcal{O}_t\mathcal{O}_t}$ are the corresponding sub-blocks. Define the four-dimensional sorting key as $z_t^{(i)} = \operatorname{vec}\left( \mu_{t\mid t-1}^{\mathcal{O}_t,(i)} \right) \in \mathbb{R}^4$. Let $\psi_t:\mathbb{R}^4\rightarrow[0,1]^4$ standardise and map these coordinates to the unit hypercube, and let $h:[0,1]^4\rightarrow[0,1]$ denote the inverse Hilbert-curve mapping. Find a permutation $\sigma_t$ such that
\begin{equation}
    h\left(\psi_t(z_t^{(\sigma_t(1))})\right) \leq \cdots \leq h\left(\psi_t(z_t^{(\sigma_t(N))})\right). \nonumber
\end{equation}
Use the previous normalised particle weights in the same order,
\begin{equation}
    \left( W_{t-1}^{(\sigma_t(1))}, \ldots, W_{t-1}^{(\sigma_t(N))} \right). \nonumber
\end{equation}

Once an ancestor is selected, its complete Rao-Blackwellised component, including the means of the non-playing teams, is retained and propagated.

% For the standard SQMC justification, sort the full Rao-Blackwellised particle state, which is the full component mean $\mu_{t-1}^{(i)} \in \mathbb{R}^{2K}$. Let $\psi: \mathbb{R}^{2K} \to [0,1]^{2K}$ standardize and quantize the state, and let $h(\cdot)$ denote the Hilbert-curve index. Find the permutation $\sigma_t$ that sorts the previous $N$ particles by their Hilbert index,
% \begin{equation}
%     h\big(\psi(\mu_{t-1}^{(\sigma_t(1))})\big) \le h\big(\psi(\mu_{t-1}^{(\sigma_t(2))})\big) \le \cdots \le h\big(\psi(\mu_{t-1}^{(\sigma_t(N))})\big), \nonumber
% \end{equation}
% and use the previous normalized weights in sorted order $W_{t-1}^{(\sigma_t(1:N))}$. We use a four-dimensional projected Hilbert ordering based on the teams involved in the current match. 

% This does not satisfy the exact transition-sufficiency condition used by Chopin and Gerber (2017), because non-playing component means remain relevant to later transitions. It is therefore treated as a dimension-reduction heuristic.

\textbf{3 Resampling by inverse transform}: Sort the complete QMC pairs $(u_t^{(i)}, v_t^{(i)})$ by their first coordinate so that $u_t^{(1)} \le \cdots \le u_t^{(N)}$, preserving the pairing between each $u_t^{(i)}$ and $v_t^{(i)}$. Perform ancestor selection using the Hilbert-ordered weights $W_{t-1}$ and the first coordinate $u_t$,
\begin{equation}
    a_i = \min \left\{ j : W^{(j)} \ge u_t^{(i)} \right\}, \qquad W^{(j)} = \sum_{k=1}^{j} W_{t-1}^{(\sigma_t(k))}, \qquad b_i = \sigma_t(a_i). \nonumber
\end{equation}

\textbf{4 Propagate}: Propagate each chosen ancestor $b_i$ through the OU dynamics of Section~\ref{sec:model}, using the predictive mean $\mu_{t\mid t-1}^{\mathcal{O}_t,(b_i)}$ and covariance $\Sigma_{t\mid t-1}$ defined in Step 2. The playing-team coordinates are drawn from the predictive marginal by the inverse-CDF transform,
\begin{equation}
    x_t^{\mathcal{O}_t,(i)} = \mu_{t\mid t-1}^{\mathcal{O}_t,(b_i)} + L_{t\mid t-1}\, \Phi^{-1}\left( v_t^{(i)} \right), \nonumber
\end{equation}
where $L_{t\mid t-1}$ is the Cholesky factor of the playing-team block $\Sigma_{t\mid t-1}^{\mathcal{O}_t\mathcal{O}_t} = \Gamma_{t\mid t-1}^{\mathcal{O}_t\mathcal{O}_t} \otimes B$, and $\Phi^{-1}$ is the coordinate-wise standard normal quantile, so that the uniform RQMC coordinates $v_t^{(i)}$ are transformed to the Gaussian predictive marginal. Proposition~\ref{prop:invcdf} in Appendix~\ref{app:inv-cdf-proof} shows that this transform samples the predictive marginal correctly. Recover the non-playing teams' coordinates in closed form via the OU prediction and the Kalman conditioning of Section~\ref{sec:model}.

\textbf{5 Compute log-weights}: Because the ancestor weights were already used in the inverse-transform resampling, compute the unnormalised log-weights for each particle $(i)$ using only the likelihood function $G_t$,
\begin{equation}
    \log \tilde{w}_t^{(i)} = \log G_t\left( y_t \mid x_t^{\mathcal{O}_t,(i)} \right), \nonumber
\end{equation}
and normalise,
\begin{equation}
    w_t^{(i)} = \frac{\tilde{w}_t^{(i)}}{\sum_{n=1}^N \tilde{w}_t^{(n)}}. \nonumber
\end{equation}

RB-SQMC Smoothing can be performed with the same algorithm as described in Section~\ref{sec:rb-smc-smoothing}, but with the random draws replaced by QMC points. At each backward step $t$, generate $M$ QMC points $(u_t^{1:M}, v_t^{1:M})$ of dimension $1 + 2K$, sort them by the conditioning state $X_{t+1}^{*(m)}$ to preserve low-discrepancy, and propagate analytically by selecting the Gaussian component via inverse transform of the backward weights and drawing
\begin{equation}
    X_t^{*(m)} = \mu_{t \mid t+1}^{(i)} + L_{t \mid t+1}\, \Phi^{-1}\left( v_t^{(m)} \right), \nonumber
\end{equation}
where $L_{t \mid t+1}$ is the Cholesky factor of the backward covariance $\Sigma_{t \mid t+1}$.

% \section{RB-SQMC Smoothing}
% \label{sec:rb-sqmc-smoothing}

\section{Parameter Estimation}
\label{sec:parameter-estimation}

Our static parameters for the model are $\Theta = (\Gamma_0, B, \kappa, \alpha, \beta)$ where $\Gamma_0 \in \mathbb{R}^{K \times K}$, $B \in \mathbb{R}^{2 \times 2}$, $\kappa > 0$, $\alpha \in \mathbb{R}$ and $\beta \in \mathbb{R}$, subject to $\Gamma_0 \succ 0$, $B \succ 0$ and $\det(B) = 1$. We fix the initial mean $\mu_0 = \mathbf{0}_{2K}$, so it is not included in $\Theta$. This is because of an identifiability constraint: If $\mu_0$ were allowed to vary, the model would be invariant to a common shift of all strengths combined with an opposite shift of the intercept $\alpha$ in the likelihood function Equation~\eqref{eq:likelihood}. This assumption ties naturally into our model: Given no prior information about teams, we assume that they start from a common baseline of zero strength, and the OU process mean-reverts to this baseline over time.

The smoothing algorithm requires the backward-sampled trajectories to be sufficiently diverse to provide a good estimate of the smoothing distribution. However, drawing many independent backward trajectories is computationally expensive: each trajectory is sampled sequentially from time $T$ to $0$, and at every time step $t$ the backward weights must be evaluated against all $N$ particles. Producing $M$ trajectories over $T$ time steps therefore costs $\mathcal{O}(MNT)$ particle-weight evaluations. Moreover, if the forward filter degenerates, the backward weights concentrate on only a few particles, so the effective number of distinct trajectories is small. The resulting smoothed trajectories then carry little diversity, which biases the estimate of the initial between-team covariance $\Gamma_0$: the backward paths are drawn from too few ancestors and the inferred spread is unreliable. 

A simpler and more efficient approach is to estimate the gradient of the marginal log-likelihood. By Fisher's identity \citep{delmoral2004feynman,delmoraldoucetsingh2010forwardsmoothing}, the marginal log-likelihood of the particle filter takes the form

\begin{equation}\nabla_\Theta \log p(y_{1:T} \mid \Theta) = \int p(x_{0:T} \mid y_{1:T}, \Theta) \nabla_\Theta \log p(x_{0:T}, y_{1:T} \mid \Theta) dx_{0:T}\end{equation}

and can be approximated by the score-function estimator described by \citet{poyiadjisdoucetsingh2011particlescore}; see \citet{kantasetal2015particleparameter} for a survey of particle methods for parameter estimation. Since complete ancestral trajectories are available from the forward filter, we use

\begin{equation}
    \widehat{S}_{N,T}(\Theta) = \sum_{i=1}^N w_T^{(i)} \nabla_\Theta \log p(x_{0:T}^{(i)}, y_{1:T} \mid \Theta)
\end{equation}

where $w_T^{(i)}$ are the (normalised) particle weights at the final time $T$ and $\widetilde{x}_{0:T}^{(i)}$ is the complete ancestral lineage of terminal particle $i$. This estimator is consistent under standard assumptions but generally biased for finite $N$. 

Furthermore, we can estimate the log-likelihood gradient using the stop-gradient trick described by \citet{scibior2021differentiableparticlefiltering} without modifying our filter. With resampling at every step and carried numerical weights $1/N$, the normalising-constant estimator is

\begin{equation}
    \widehat Z_0 = 1, \qquad \widehat Z_t = \widehat Z_{t-1} \left[ \frac{1}{N} \sum_{i=1}^N G_t(y_t \mid x_t^{\mathcal{O}_t,(i)}) \right], \nonumber
\end{equation}

so that
\begin{equation}
    \log \widehat Z_T = \sum_{t=1}^T \log \left[ \frac{1}{N} \sum_{i=1}^N G_t(y_t \mid x_t^{\mathcal{O}_t,(i)}) \right]. \nonumber
\end{equation}

We therefore maximise the marginal log-likelihood $\log p(y_{1:T} \mid \Theta)$ by gradient ascent on the static parameters $\Theta$. The gradient is obtained by backpropagating through the normalising-constant estimator $\log \widehat Z_T$, which is differentiable with respect to $\Theta$ because the OU transition and the bivariate-Poisson likelihood are differentiable. The stop-gradient trick of \citet{scibior2021differentiableparticlefiltering} makes the resampling step differentiable by blocking the gradient through the resampled ancestor indices, so that the resampling does not contribute a spurious gradient while the transition and likelihood terms remain differentiable.

\section{Performance comparison}
\label{sec:performance-comparison}

We attempt to model the performance of teams during the World Cup 2026. There are $K = 48$ teams competing in the World Cup 2026.

We compare RB-SQMC against a factorial state-space model, specifically a moment-based extended Kalman filter (EKF) baseline implemented in the \texttt{cuthberto-carlos} package \citep{cuthbertocarlos2026}, which builds on the \texttt{cuthbert} library for factorial filtering. We are comparing two different methods for approximating the same underlying model, RB-SQMC which uses a particle approximation of the latent state and EKF which uses a biased Gaussian approximation. 

% The comparison is meant to evaluate the performance of the two methods in terms of their ability to estimate the model parameters and make predictions.

% The comparison serves two purposes. First, it isolates the effect of the correlated Rao-Blackwellised construction developed in this chapter: the EKF is a factorial Gaussian moment filter that maintains an independent marginal for each team and discards the cross-team covariance after every match, whereas RB-SQMC retains that dependence through the Rao-Blackwell covariance update. Second, it contrasts a closed-form approximation with a particle method, so that the cost of the particle representation is made explicit. The two methods are not intended to be interchangeable; the EKF is a fast baseline whose Gaussian approximation is exact only under linear-Gaussian observation, which the bivariate-Poisson likelihood violates.

\subsection{Experimental design}

Both methods consume a single chronologically ordered data set of football results, so that the comparison is not confounded by differences in data preparation. The data are split chronologically into a training period $[1980\text{-}01\text{-}01, 2024\text{-}01\text{-}01)$, a test period $[2024\text{-}01\text{-}01, 2026\text{-}06\text{-}11)$, and a prediction period from 11 June 2026. The prediction period contains the four World Cup 2026 fixtures used for the headline scoring. Matches with scores exceeding the truncation threshold $\max\_goals = 8$ are excluded before the split, and same-day repeats are retained as separate rows with zero elapsed time.

The static parameters $\Theta = (\Gamma_0, B, \kappa, \alpha, \beta)$ are estimated by maximising the marginal log-likelihood as mentioned in Section~\ref{sec:parameter-estimation}. Both methods use Adam with a cosine-decay learning-rate schedule, learning rate $0.02$, seed $0$ and exactly $50$ updates. Because the SQMC log-likelihood is a particle estimate, its gradient is averaged over $10$ independently randomised RQMC replicas at each update, whereas the EKF gradient is deterministic. The EKF integrates the bivariate-Poisson likelihood under the predictive Gaussian by Gauss--Hermite quadrature of degree $32$ \citep{ghq}; RB-SQMC evaluates the same likelihood on the sampled playing-team coordinates. The full parameter set and the exact configuration are listed in Appendix~\ref{app:config}.

\subsection{Training: marginal log-likelihood}

Table~\ref{tab:logz} reports the training and test log-likelihoods. 

% The two columns are not directly comparable as exact marginal likelihoods: the EKF value is a Gaussian-approximate normalising constant from moment filtering, while the SQMC value is a particle estimate with finite-sample error. The two methods also adopt different time-step conventions, the EKF propagating only the two playing teams by each team's last-appearance gap and RB-SQMC propagating the full covariance per row. The raw gap must therefore not be read as a superiority verdict.

\begin{table}[htbp]
\centering
\begin{tabular}{lcc}
\toprule
 & EKF & RB-SQMC \\
\midrule
Final train log-likelihood & $-1.568\times10^{4}$ & $-1.499\times10^{4}$ \\
Final test log-likelihood & $-1.040\times10^{3}$ & $-1.005\times10^{3}$ \\
Best test log-likelihood (epoch) & $-1.040\times10^{3}$ (50) & $-9.709\times10^{2}$ (26) \\
\bottomrule
\end{tabular}
\caption{Training and test log-likelihoods for the EKF and RB-SQMC. The EKF value is a Gaussian-approximate normalising constant; the SQMC value is a particle estimate. The two are not directly comparable as exact marginal likelihoods.}
\label{tab:logz}
\end{table}

Figure~\ref{fig:logz-overlay} overlays the training and test log-likelihood trajectories across epochs for both methods.

\begin{figure}[htbp]
\centering
\includegraphics[width=0.9\textwidth]{results/sqmc_ekf/combined/images/logz_overlay_train_test.png}
\caption{Training and test log-likelihood against epoch for the EKF and RB-SQMC. The EKF value is a Gaussian-approximate normalising constant; the SQMC value is a particle estimate.}
\label{fig:logz-overlay}
\end{figure}

Within the range examined, RB-SQMC attained a higher final train log-likelihood than the EKF ($-1.499\times10^{4}$ versus $-1.568\times10^{4}$) and a higher best test log-likelihood ($-9.709\times10^{2}$ at epoch 26 versus $-1.040\times10^{3}$ at epoch 50). This is consistent with the particle method evaluating the non-Gaussian likelihood without the Gaussian approximation. The comparison is nevertheless confounded in two ways. First, the EKF value is a Gaussian-approximate normalising constant from moment filtering, whereas the SQMC value is a particle estimate with finite-sample error, so the raw gap is not a direct measure of model fit. Second, the two methods adopt different time-step conventions, the EKF propagating only the two playing teams by each team's last-appearance gap and RB-SQMC propagating the full covariance per row. The raw gap must therefore not be read as a superiority verdict.

The SQMC test log-likelihood is also markedly less stable across epochs than the EKF's. The EKF trajectory is smooth and monotone, improving steadily from $-1.102\times10^{3}$ to $-1.040\times10^{3}$ over the 50 updates. The SQMC trajectory oscillates between approximately $-9.7\times10^{2}$ and $-1.06\times10^{3}$, with the best value reached at epoch 26 rather than at the final checkpoint. This instability reflects the particle estimate's Monte Carlo error, which the ten-replica gradient average reduces but does not eliminate; the best test value is therefore a favourable draw in a noisy trajectory rather than a robust optimum.

\subsection{Prediction}

For each prediction-split match, both methods produce a one-step-ahead score grid before assimilating the known result, so that the current score cannot enter its own forecast. We score the resulting distributions with the mean Brier score over the three outcomes (home win, draw, away win), the mean predictive log-score, and outcome accuracy, all restricted to the four World Cup 2026 fixtures. The Brier score is defined as,

\begin{equation}
    \text{Brier score} = \frac{1}{N} \sum_{i=1}^N \sum_{j=1}^{N} (p_{ij} - o_{ij})^2, \nonumber
\end{equation}

where $p_{ij}$ is the predicted probability of outcome $j$ for match $i$, and $o_{ij}$ is the observed outcome indicator. The predictive log-score is defined as,

\begin{equation}
    \text{Predictive log-score} = \frac{1}{N} \sum_{i=1}^N \log p_{i,o_i}, \nonumber
\end{equation}

which is the average log-probability assigned to the observed outcome $o_i$ for match $i$. Outcome accuracy is defined as the proportion of matches for which the predicted outcome with the highest probability matches the observed outcome.

Table~\ref{tab:prediction} reports the results, and the full win--loss predictions for the fixtures and the knockout stage are collected in Appendix~\ref{app:prediction}.

\begin{table}[htbp]
\centering
\begin{tabular}{lcc}
\toprule
 & EKF & RB-SQMC \\
\midrule
Mean Brier score & $0.5442$ & $0.6521$ \\
Uniform reference Brier score & $2/3$ & $2/3$ \\
Brier skill score vs.\ uniform & $0.1838$ & $0.0218$ \\
Mean predictive log-score & $-3.131$ & $-3.523$ \\
Outcome accuracy & $0.5288$ & $0.5096$ \\
Exact-score accuracy & $0.0769$ & $0.0673$ \\
Scored matches & $104$ & $104$ \\
\bottomrule
\end{tabular}
\caption{One-step-ahead prediction metrics on the World Cup 2026 fixtures. The Brier skill score is relative to the uniform reference $2/3$.}
\label{tab:prediction}
\end{table}

Table~\ref{tab:pre-tournament} reports the pre-tournament top-ten team rankings under each method, ordered by total strength (the sum of the attacking and defensive components) before the World Cup.

\begin{table}[htbp]
\centering
\begin{tabular}{clll}
\toprule
Rank & EKF & RB-SQMC \\
\midrule
1 & Belgium & Brazil \\
2 & Brazil & England \\
3 & Switzerland & Portugal \\
4 & Colombia & Belgium \\
5 & Austria & France \\
6 & Portugal & Japan \\
7 & France & Sweden \\
8 & Algeria & Colombia \\
9 & Ecuador & Germany \\
10 & Japan & Argentina \\
\bottomrule
\end{tabular}
\caption{Pre-tournament top-ten team rankings under the EKF and RB-SQMC, ordered by total strength before the World Cup.}
\label{tab:pre-tournament}
\end{table}

We review the final match of the tournament, Spain vs.\ Argentina, played on 19 July 2026. The images below show the predicted score grids for both methods, with the actual result (a $1$--$0$ win for Spain) marked by a red cross. The EKF assigns the highest probability to a draw ($0.207$) and a home win ($0.479$), whereas RB-SQMC assigns the highest probability to a home win ($0.537$); both methods place the modal score at $1$--$1$ or $2$--$1$, neither of which matches the actual $1$--$0$ result.

\begin{figure}[htbp]
\centering
\begin{subfigure}{0.9\textwidth}
\includegraphics[width=\textwidth]{results/sqmc_ekf/combined/images/ekf_predictions/prediction_plots/2026-07-19_Spain_vs_Argentina.png}
\caption{EKF predicted score grid.}
\label{fig:prediction-final-ekf}
\end{subfigure}
\begin{subfigure}{0.9\textwidth}
\includegraphics[width=\textwidth]{results/sqmc_ekf/combined/images/sqmc_predictions/prediction_plots/2026-07-19_Spain_vs_Argentina.png}
\caption{RB-SQMC predicted score grid.}
\label{fig:prediction-final-sqmc}
\end{subfigure}
\caption{Predicted score grids for the final match, Spain vs.\ Argentina, under the EKF and RB-SQMC. The actual result is marked by a red cross.}
\label{fig:prediction-final}
\end{figure}

% Table~\ref{tab:prediction-knockout} reports the win--loss predictions of the bivariate-Poisson model for the two semi-finals and the final, alongside the actual results.

% \begin{table}[htbp]
% \centering
% \begin{tabular}{llccc}
% \toprule
% Home team & Away team & Actual result & EKF predicted & RB-SQMC predicted \\
% \midrule
% England & Argentina & \placeholder{result} & \placeholder{outcome} & \placeholder{outcome} \\
% France & Spain & \placeholder{result} & \placeholder{outcome} & \placeholder{outcome} \\
% Spain & Argentina & \placeholder{result} & \placeholder{outcome} & \placeholder{outcome} \\
% \bottomrule
% \end{tabular}
% \caption{Win--loss predictions of the bivariate-Poisson model for the two semi-finals and the final. The predicted outcome is the mode of the three-outcome distribution under each method.}
% \label{tab:prediction-knockout}
% \end{table}

On the 104 prediction-split matches, the EKF produced the lower mean Brier score ($0.5442$ versus $0.6521$) and the higher Brier skill score relative to the uniform reference ($0.1838$ versus $0.0218$). The EKF also attained the higher mean predictive log-score ($-3.131$ versus $-3.523$) and the higher outcome accuracy ($0.5288$ versus $0.5096$). RB-SQMC is therefore barely better than the uniform reference on the Brier skill score, whereas the EKF is substantially better. This is a negative result for RB-SQMC within the range examined, and it is reported directly. It should not be over-interpreted: 104 matches cannot fully separate the methods, the SQMC forecast is a single particle realisation with no averaging over scrambles, and the EKF's advantage may reflect the particular fixtures rather than a systematic property. The appropriate reading is that the prediction protocol runs end to end and produces well-formed, properly normalised score grids for both methods; a definitive performance comparison would require more held-out matches and averaging over scrambles.

\subsection{Evaluation}

Figure~\ref{fig:timeseries} shows the filtered attacking and defensive strength trajectories for a representative team under each method. The EKF produces a deterministic linear trajectory, whereas RB-SQMC reports the particle-weighted posterior mean, $\bar{x}_t^{(k)} = \sum_{i=1}^N w_t^{(i)} x_t^{(i,k)}$, where $w_t^{(i)}$ are the normalised particle weights after the match at time $t$ has been assimilated. The figure plots this weighted mean only; the particle spread is not shown, so the RB-SQMC trajectories should be read as posterior mean estimates rather than as means with uncertainty bands.

\begin{figure}[htbp]
\centering
\begin{subfigure}{0.9\textwidth}
\includegraphics[width=\textwidth]{results/sqmc_ekf/combined/images/ekf_timeseries_states.png}
\caption{EKF filtered states.}
\label{fig:timeseries-ekf}
\end{subfigure}
\begin{subfigure}{0.9\textwidth}
\includegraphics[width=\textwidth]{results/sqmc_ekf/combined/images/sqmc_timeseries_states.png}
\caption{RB-SQMC filtered states.}
\label{fig:timeseries-sqmc}
\end{subfigure}
\caption{Filtered attacking and defensive strengths over time for a representative team under the EKF and RB-SQMC. The EKF trajectories are deterministic moment estimates; the RB-SQMC trajectories are particle-weighted posterior means.}
\label{fig:timeseries}
\end{figure}

Figure~\ref{fig:topn-correlation} reports the correlation between the RB-SQMC top-$N$ rankings before and after the World Cup, which summarises how much the tournament results reorder the strongest teams.

\begin{figure}[htbp]
\centering
\includegraphics[width=0.9\textwidth]{results/sqmc_ekf/combined/images/sqmc_correlation_topn_bar.png}
\caption{Correlation between the RB-SQMC top-$N$ team rankings before and after the World Cup, as a function of $N$.}
\label{fig:topn-correlation}
\end{figure}

\subsection{Comparison with the FIFA ranking}

The quantitative metrics above favour the EKF on prediction but RB-SQMC on the log-likelihood. To assess which method produces the more plausible ranking of team strengths, we compare the post-tournament rankings of both methods with the FIFA Men's World Ranking \citep{fifaworldranking2026}. Since August 2018 the FIFA ranking has been computed by an Elo-style rating system, in which each team's points are updated after every match according to the expected-result formula
\begin{equation}
P_{\text{after}} = P_{\text{before}} + I\,(W - W_e), \qquad W_e = \frac{1}{10^{-\Delta/c} + 1},
\end{equation}
where $I$ is a match-importance coefficient, $W$ is the match result, $\Delta$ is the pre-match rating difference and $c = 600$ is a scale constant. The FIFA ranking therefore provides an independent, continuously updated measure of relative team strength against which the two model rankings can be judged.

Table~\ref{tab:fifa-comparison} reports the top ten teams under each method after the World Cup, together with the FIFA ranking published on 20 July 2026.

\begin{table}[htbp]
\centering
\begin{tabular}{clll}
\toprule
Rank & EKF & RB-SQMC & FIFA \\
\midrule
1 & Spain & England & Spain \\
2 & Argentina & France & Argentina \\
3 & Brazil & Portugal & France \\
4 & Belgium & Scotland & England \\
5 & Portugal & Germany & Brazil \\
6 & Colombia & United States & Morocco \\
7 & England & Japan & Portugal \\
8 & Morocco & Belgium & Belgium \\
9 & Switzerland & Argentina & Netherlands \\
10 & France & Brazil & Mexico \\
\bottomrule
\end{tabular}
\caption{Post-tournament top-ten team rankings under the EKF and RB-SQMC, compared with the FIFA Men's World Ranking of 20 July 2026 \citep{fifaworldranking2026}.}
\label{tab:fifa-comparison}
\end{table}

The EKF's post-tournament ranking aligns closely with the FIFA ranking: both place Spain first and Argentina second, and both include France, England, Brazil, Morocco, Portugal and Belgium in the top ten. The EKF's principal discrepancies are the omission of the Netherlands and Mexico and the inclusion of Colombia and Switzerland. RB-SQMC's ranking is less plausible: it omits Spain from the top ten entirely, places Scotland fourth and the United States sixth, and ranks Argentina ninth despite the FIFA ranking placing it second. The EKF therefore produces the more justifiable qualitative ranking, consistent with its better prediction metrics, even though RB-SQMC attains the higher log-likelihood. The full pre- and post-tournament rankings for both methods are collected in Appendix~\ref{app:rankings}.

\subsection{Discussion and Limitations}

The comparison is designed so that any difference between the methods is attributable to the correlated Rao-Blackwellised construction rather than to data preparation, because both methods consume the same frozen chronological data set and the same optimiser. Within the range examined, RB-SQMC attained a higher test log-likelihood, but the EKF produced the better prediction scores on the World Cup fixtures and the more plausible ranking of team strengths. These observations are not contradictory: the log-likelihood comparison is confounded by the approximation status of the EKF normalising constant and by the different time-step conventions, while the prediction comparison is based on a limited number of held-out matches.

The computational cost of RB-SQMC is a material limitation. The EKF is a closed-form moment filter: it compiled in $0.97$\,s and executed in $12.5$\,s, with prediction in $0.96$\,s. RB-SQMC compiled in $16.0$\,s and executed in $3.66\times10^{3}$\,s (approximately one hour) on an A100 GPU, with prediction in $17.5$\,s, dominated by the repeated particle-filter evaluations and the Hilbert sorting. The comparison therefore separates a statistical question from a computational one: RB-SQMC is a particle method whose cost scales with the particle count and the number of optimisation epochs, while the EKF is a fast approximation. For a sequential forecasting problem in which match results arrive at low frequency, the RB-SQMC filtering cost is acceptable, but the parameter-estimation cost is the material bottleneck. A more efficient estimator of the static parameters, for example a variational or stochastic-gradient procedure, would reduce this start-up cost.

The SQMC test log-likelihood's instability across epochs, reported in Section~\ref{sec:performance-comparison}, indicates that the particle estimate is not fully converged at $128$ particles and $10$ gradient replicas. Increasing the particle count or the number of replicas would reduce the Monte Carlo error, at a further computational cost. The prediction metrics are also restricted to a single particle realisation with no averaging over scrambles; averaging over several RQMC scrambles would give a more reliable estimate of RB-SQMC's predictive performance.

The model assigns low probability to extreme scorelines. In the knockout stage, the France--England match ended $4$--$6$, a result to which both methods assigned very low probability (EKF log-likelihood $-8.43$, RB-SQMC $-11.29$). The bivariate-Poisson likelihood with a shared component does not capture the tail behaviour of very high-scoring matches, and the truncation at $\max\_goals = 8$ excludes the most extreme outcomes from the training data. A more flexible observation model, or a heavier-tailed score distribution, would be required to estimate these probabilities reliably.

Although the model is developed for the World Cup 2026, the same construction applies to league matches: the summer break between seasons also involves other games, so the sparse pairwise likelihood and the Rao-Blackwellised sampling of the playing teams carry over directly. A more exact comparison of the two methods would compare the predicted probabilities with those implied by prediction markets (for example Polymarket, Kalshi or sports bookmakers), which aggregate the beliefs of many market participants and provide an external benchmark for calibration.